documentclass[
    reprint,
    aps,
    pra,
    amsmath,
    amssymb,
    superscriptaddress,
    onecolumn,
    11pt
]{revtex4-2}

\usepackage{graphicx}
\usepackage{dcolumn}
\usepackage{bm}
\usepackage{hyperref}
\usepackage{xcolor}

\hypersetup{
    colorlinks=true,
    linkcolor=blue!50!black,
    citecolor=blue!50!black,
    urlcolor=blue!50!black
}

\begin{document}

\title{Quick Research Note: [Short Descriptive Title]}

\author{[Author Name]}
\altaffiliation[Also at: ]{[Department/Laboratory Name]}
\affiliation{%
 [Research Group or Laboratory Name], [University or Institution Name], [City, State, Zip Code]
}%

\date{\today}

\begin{abstract}
  This document serves as a standardized template for quick technical write-ups, theoretical derivations, or experimental progress reports within the research group. Briefly summarize the core objective, key methodology, principal findings, and next steps in 3-5 sentences. Here is a citation of paper \cite{johndoe2024}.
\end{abstract}

\maketitle

\section{\label{sec:intro}Introduction and Motivation}
State the context of the current note. Why is this calculation, derivation, or diagnostic being performed? Connect it to ongoing group projects or recent experimental/theoretical roadblocks.

\section{\label{sec:formulation}Theoretical Formulation / Setup}
Define the mathematical framework or physical model being used. For example, consider a Hamiltonian of the general form:
\begin{equation}
    \hat{H} = \sum_{i} \epsilon_i \hat{c}^\dagger_i \hat{c}_i + \sum_{ij} t_{ij} \left( \hat{c}^\dagger_i \hat{c}_j + \text{h.c.} \right) + U \sum_i \hat{n}_{i\uparrow} \hat{n}_{i\downarrow},
    \label{eq:hamiltonian}
\end{equation}
where the symbols retain their standard definitions within the context of our group's convention.

\section{\label{sec:results}Calculations and Discussion}
Describe the derivation steps or computational results. When referencing equations, use Eq.~(\ref{eq:hamiltonian}).

\subsection{Analytical Limits}
In the limit of strong interactions ($U \gg t$), we can perform a perturbative expansion to second order in $t/U$, yielding an effective superexchange interaction:
\begin{equation}
    J = \frac{4t^2}{U}.
\end{equation}
This confirms our initial hypothesis regarding the low-energy effective spin model.

\section{\label{sec:conclusion}Summary and Action Items}
Conclude with clear takeaways and what needs to be tested next.
\begin{itemize}
    \item \textbf{Finding 1:} The perturbative limit maps correctly onto the Heisenberg model.
    \item \textbf{Action Item 1:} Extend the numerical code to include next-nearest-neighbor hopping terms by [Date].
    \item \textbf{Action Item 2:} Discuss discrepancies with experimental dataset B during the next group meeting.
\end{itemize}

\newpage
\appendix

\section{Derivation Details}
\subsection{Second-Order Perturbation Theory}
Here we include lengthy algebraic manipulations or supplementary proofs that would otherwise interrupt the flow of the main text.
\begin{equation}
    E_0^{(2)} = \sum_{m \neq 0} \frac{|\langle m | \hat{H}_1 | 0 \rangle|^2}{E_0^{(0)} - E_m^{(0)}}.
\end{equation}

\newpage
\bibliography{reference}

\end{document}
